\documentclass[14pt]{extarticle}
\usepackage{natbib}

\usepackage{xcolor}
\usepackage{amsmath}
\newcommand{\tuple}[1]{ \langle #1 \rangle }
%\usepackage{automata}
\usepackage{times}
\usepackage{ltablex}
\usepackage{tasks}
\usepackage{threeparttable}
\usepackage{booktabs}
\usepackage{xcolor} % For custom colors
\usepackage{tikz} % For styling enumerate numbers
\usepackage{tcolorbox} % For colored box styling
\usepackage{amsmath, amsfonts, amssymb, amsthm} % Math related
\usepackage{natbib}
\usepackage{fontspec}
\usepackage{luatexja}
\usepackage[mathscr]{euscript}

%%%%%% Template
\usepackage{hyperref}
\hypersetup{colorlinks=true,allcolors=blue}

\usepackage{vmargin}
\setpapersize{USletter}
\setmarginsrb{1.0in}{1.0in}{1.0in}{0.6in}{0pt}{0pt}{0pt}{0.4in}

% HOW TO USE THE ABOVE:
%\setmarginsrb{leftmargin}{topmargin}{rightmargin}{bottommargin}{headheight}{headsep}{footheight}{footskip}
%\raggedbottom
% paragraphs indent & skip:
\parindent  0.3cm
\parskip    -0.01cm

\usepackage{tikz}
\usetikzlibrary{backgrounds}

% hyphenation:
% \hyphenpenalty=10000 % no hyphen
% \exhyphenpenalty=10000 % no hyphen
\sloppy

% notes-style paragraph spacing and indentation:
\usepackage{parskip}
\setlength{\parindent}{0cm}

% let derivations break across pages
\allowdisplaybreaks

\newcommand{\orange}[1]{\textcolor{orange}{#1}}
\newcommand{\blue}[1]{\textcolor{blue}{#1}}
\newcommand{\red}[1]{\textcolor{red}{#1}}
\newcommand{\freq}[1]{{\bf \sf F}(#1)}
\newcommand{\datafreq}[2]{{{\bf \sf F}_{#1}(#2)}}

\def\qqquad{\quad\qquad}
\def\qqqquad{\qquad\qquad}

%%%%%%%%%%%%%%%%%%%%%%%%%%%%%%%%%%%%%%%%%%%%%%%%%%%%%%%%%%%%%%%%%%%%%%%%%%%%%%%%
%%%%%%%%%%%%%%%%%%%%%%%%%%%%%%%%%%%%%%%%%%%%%%%%%%%%%%%%%%%%%%%%%%%%%%%%%%%%%%%%

% fill-in-blank question style, found in https://tex.stackexchange.com/a/505089

\usepackage{ifthen}
\usepackage{tocloft}
\usepackage{exercise}
% \usepackage{xcolor}

% Set the Show Answers Boolean
\newboolean{showAns}
\setboolean{showAns}{false}
\newcommand{\showAns}{\setboolean{showAns}{true}}

% The length of the Answer line
\newlength{\answerlength}
\newcommand{\anslen}[1]{\settowidth{\answerlength}{#1}}

% ans command that indicates space for an answer or shows the answer in red
\newcommand{\ans}[1]{\settowidth{\answerlength}{\hspace{2ex}#1\hspace{2ex}}%
    \ifthenelse{\boolean{showAns}}%
        {\textcolor{red}{\underline{\hspace{2ex}#1\hspace{2ex}}}}%
        {\underline{\hspace{\answerlength}}}}%

\newcommand{\details}[1]{\settowidth{\answerlength}{#1}%
    \ifthenelse{\boolean{showAns}}%
        {\begin{tcolorbox}[colback=blue!5!white,colframe=blue!75!black,title=Solution] #1 \end{tcolorbox}}%
        {}}%

% Formatting how multiple choices Questions are formated.
\settasks{label=(\Alph*), label-width=30pt}


% Some commands for the Exercise Question package
\renewcommand{\QuestionNB}{\Large\protect\textcircled{\small\bfseries\arabic{Question}}\ }
\renewcommand{\ExerciseHeader}{} %no header
\renewcommand{\QuestionBefore}{3ex} %Space above each Q
\setlength{\QuestionIndent}{8pt} % Indent after Q number


% To create the list of answers with tocloft...
\newcommand{\listanswername}{Answers}
\newlistof[Question]{answer}{Answers}{\listanswername}

% Creates a TOC for Answers
\newcounter{prevQ}
\newcommand{\answer}[1]{\refstepcounter{answer}%
\ans{#1}%
\ifnum\theQuestion=\theprevQ%
        \addcontentsline{Answers}{answer}{\protect\numberline{}#1}% don't include the Q number
        \else%
        \addcontentsline{Answers}{answer}{\protect\numberline{\theQuestion}#1}%
        \setcounter{prevQ}{\value{Question}}%
        \fi%
        }%

% \hyphenpenalty=10000 % no hyphen
% \exhyphenpenalty=10000 % no hyphen
\sloppy              % hyphen

\newcommand{\HRule}{\rule{\linewidth}{0.5mm}}
\newcommand{\Hrule}{\rule{\linewidth}{0.3mm}}

%tocloft formatting listofanswers
\renewcommand{\cftAnswerstitlefont}{\bfseries\large}
\renewcommand{\cftanswerdotsep}{\cftnodots}
\cftpagenumbersoff{answer}
\addtolength{\cftanswernumwidth}{10pt}

\makeatletter% since there's an at-sign (@) in the command name
\renewcommand{\@maketitle}{%
  \parindent=0pt% don't indent paragraphs in the title block
  \centering
  {\Large \bfseries\textsc{\@title}} \\
  \vspace{5pt}
  {\large \textit{\@author}} \\
  \HRule \\
  \vspace{1em}
}
\makeatother% resets the meaning of the at-sign (@)


% --- Fonts (robust fallback) ---
\IfFontExistsTF{Crimson Pro Light}{
  \setmainfont{Crimson Pro Light}[
    ItalicFont={* Italic},
    BoldFont={Crimson Pro Medium},
    BoldItalicFont={Crimson Pro Medium Italic}]
  \setsansfont{Crimson Pro Light}[
    ItalicFont={* Italic},
    BoldFont={Crimson Pro Medium},
    BoldItalicFont={Crimson Pro Medium Italic}]
}{
  \setmainfont{HaranoAjiMincho}
  \setsansfont{HaranoAjiGothic}
}
\title{Problem Set 3: Two-Period New Keynesian Model}
\author{Macroeconomics II (Spring 2026) \\ Hui-Jun Chen}



% ------------------------------------------------------------
% Explanations for each option (only printed when \showAns is enabled)
% Usage: \ExplainChoices{A}{B}{C}{D}
% ------------------------------------------------------------
\newcommand{\ExplainChoices}[4]{%
\ifthenelse{\boolean{showAns}}{%
\vspace{0.3em}\noindent\textcolor{red}{\textbf{Explanations.}}
\begin{itemize}
\item[(A)] #1
\item[(B)] #2
\item[(C)] #3
\item[(D)] #4
\end{itemize}
}{}%
}

\begin{document}

\maketitle

% Uncomment to show answers + explanations:
\showAns

\begin{Exercise}

\section*{Problem 1 (Questions 1--20): IS / Euler block}
\Question \textbf{Timeline in two-period NK. Which expectation operator is nontrivial?}
\answer{B}
\begin{tasks}(1)
\task $\mathbb{E}_1[\cdot]$ about period 2 variables
\task $\mathbb{E}_0[\cdot]$ about period 1 variables
\task $\mathbb{E}_{-1}[\cdot]$ about period 0 variables
\task No expectations appear in a two-period model
\end{tasks}
\ExplainChoices{Wrong: no period 2 exists.}{Correct: only expectations from 0 to 1 matter.}{Wrong: not relevant here.}{Wrong: expectations can matter for period 1.}

\Question \textbf{Fisher relation (ex-ante real rate). Which is correct?}
\answer{C}
\begin{tasks}(1)
\task $r_0=i_0+\mathbb{E}_0\pi_1$
\task $r_0=\mathbb{E}_0\pi_1-i_0$
\task $r_0=i_0-\mathbb{E}_0\pi_1$
\task $r_0=i_0-\pi_0$
\end{tasks}
\ExplainChoices{Wrong sign.}{Wrong sign.}{Correct.}{Wrong: uses current instead of expected inflation.}

\Question \textbf{NK IS (Euler) equation: pick the standard form.}
\answer{A}
\begin{tasks}(1)
\task $x_0=\mathbb{E}_0 x_1-\frac{1}{\sigma}(i_0-\mathbb{E}_0\pi_1-r_0^n)$
\task $x_0=\mathbb{E}_0 x_1+\frac{1}{\sigma}(i_0-\mathbb{E}_0\pi_1-r_0^n)$
\task $x_0=\frac{1}{\sigma}(i_0-\mathbb{E}_0\pi_1-r_0^n)$
\task $x_0=\mathbb{E}_0 x_1-\sigma(i_0-\mathbb{E}_0\pi_1-r_0^n)$
\end{tasks}
\ExplainChoices{Correct.}{Wrong sign.}{Missing expectation.}{Wrong coefficient.}

\Question \textbf{Meaning of $\sigma$ in NK IS.}
\answer{C}
\begin{tasks}(1)
\task Policy response to inflation
\task Price rigidity parameter
\task Inverse intertemporal elasticity of substitution
\task Slope of Phillips curve
\end{tasks}
\ExplainChoices{Wrong.}{Wrong.}{Correct.}{Wrong.}

\Question \textbf{If $x_1=0$ for sure, NK IS implies:}
\answer{B}
\begin{tasks}(1)
\task $x_0=\mathbb{E}_0 x_1$
\task $x_0=-\frac{1}{\sigma}(i_0-\mathbb{E}_0\pi_1-r_0^n)$
\task $x_0=+\frac{1}{\sigma}(i_0-\mathbb{E}_0\pi_1-r_0^n)$
\task $x_0=0$ always
\end{tasks}
\ExplainChoices{Wrong.}{Correct.}{Wrong sign.}{Wrong.}

\Question \textbf{Numerical IS. $\sigma=2$, $r_0^n=1\%$, $i_0=3\%$, $\mathbb{E}_0\pi_1=1\%$, $x_1=0$. Then $x_0=$}
\answer{C}
\begin{tasks}(1)
\task $+0.5\%$
\task $0$
\task $-0.5\%$
\task $-1\%$
\end{tasks}
\ExplainChoices{Wrong sign.}{Wrong.}{Correct.}{Wrong: forgot $1/\sigma$.}

\Question \textbf{Natural real rate $r_0^n$ is:}
\answer{D}
\begin{tasks}(1)
\task The nominal policy rate
\task Inflation target
\task Real rate consistent with $x_0=0$ under sticky prices
\task Real rate consistent with $x_0=0$ under flexible prices
\end{tasks}
\ExplainChoices{Wrong.}{Wrong.}{Close but not definition.}{Correct.}

\Question \textbf{Holding $i_0$ fixed, higher $\mathbb{E}_0\pi_1$ implies:}
\answer{B}
\begin{tasks}(1)
\task Higher $r_0$ and lower $x_0$
\task Lower $r_0$ and higher $x_0$
\task Lower $r_0$ and lower $x_0$
\task No effect on $x_0$
\end{tasks}
\ExplainChoices{Wrong.}{Correct.}{Wrong.}{Wrong.}

\Question \textbf{ZLB constraint is:}
\answer{C}
\begin{tasks}(1)
\task $r_0\ge0$
\task $\pi_0\ge0$
\task $i_0\ge0$
\task $x_0\ge0$
\end{tasks}
\ExplainChoices{Wrong.}{Wrong.}{Correct.}{Wrong.}

\Question \textbf{At ZLB, if expected inflation falls, then:}
\answer{A}
\begin{tasks}(1)
\task Real rate rises and demand falls
\task Real rate falls and demand rises
\task Nominal rate rises automatically
\task NKPC shifts down only
\end{tasks}
\ExplainChoices{Correct.}{Wrong.}{Wrong.}{Wrong.}

\Question \textbf{NK Phillips curve (NKPC) form is:}
\answer{B}
\begin{tasks}(1)
\task $\pi_0=\pi_{-1}+\kappa x_0+u_0$
\task $\pi_0=\beta\mathbb{E}_0\pi_1+\kappa x_0+u_0$
\task $\pi_0=\beta\pi_1-\kappa x_0+u_0$
\task $\pi_0=\kappa x_0+u_0$
\end{tasks}
\ExplainChoices{Backward-looking.}{Correct.}{Wrong.}{Missing expectations.}

\Question \textbf{Interpretation of $\beta$ in NKPC:}
\answer{A}
\begin{tasks}(1)
\task Discount factor weight on expected inflation
\task Policy response to inflation
\task IS slope parameter
\task Cost-push shock size
\end{tasks}
\ExplainChoices{Correct.}{Wrong.}{Wrong.}{Wrong.}

\Question \textbf{Interpretation of $u_0$ in NKPC:}
\answer{A}
\begin{tasks}(1)
\task Cost-push/markup shock
\task Demand shock
\task Natural rate
\task Nominal rate
\end{tasks}
\ExplainChoices{Correct.}{Wrong.}{Wrong.}{Wrong.}

\Question \textbf{If $\mathbb{E}_0\pi_1=0$, NKPC becomes:}
\answer{A}
\begin{tasks}(1)
\task $\pi_0=\kappa x_0+u_0$
\task $\pi_0=\beta+\kappa x_0+u_0$
\task $\pi_0=\beta\pi_0+\kappa x_0+u_0$
\task $\pi_0=u_0$
\end{tasks}
\ExplainChoices{Correct.}{Wrong.}{Wrong.}{Wrong.}

\Question \textbf{Taylor rule (simple) is:}
\answer{A}
\begin{tasks}(1)
\task $i_0=\bar i+\phi_\pi\pi_0+\phi_x x_0$
\task $i_0=r_0^n$
\task $i_0=\pi_0-\phi_\pi x_0$
\task $i_0=\mathbb{E}_0\pi_1$
\end{tasks}
\ExplainChoices{Correct.}{Wrong.}{Wrong.}{Wrong.}

\Question \textbf{Taylor principle typically requires:}
\answer{C}
\begin{tasks}(1)
\task $\phi_\pi<1$
\task $\phi_\pi=1$
\task $\phi_\pi>1$
\task $\phi_\pi<0$
\end{tasks}
\ExplainChoices{Wrong.}{Boundary.}{Correct.}{Wrong.}

\Question \textbf{Determinacy intuition: when inflation rises, real rate should:}
\answer{A}
\begin{tasks}(1)
\task Rise (tighten)
\task Fall (loosen)
\task Stay constant
\task Be undefined
\end{tasks}
\ExplainChoices{Correct.}{Wrong.}{Wrong.}{Wrong.}

\Question \textbf{Demand shock (IS) tends to move:}
\answer{A}
\begin{tasks}(1)
\task $x\uparrow,\pi\uparrow$
\task $x\downarrow,\pi\uparrow$
\task $x\uparrow,\pi\downarrow$
\task $x\downarrow,\pi\downarrow$
\end{tasks}
\ExplainChoices{Correct.}{Cost-push.}{Wrong.}{Wrong.}

\Question \textbf{Cost-push shock ($u_0>0$) tends to move:}
\answer{B}
\begin{tasks}(1)
\task $x\uparrow,\pi\uparrow$
\task $x\downarrow,\pi\uparrow$
\task $x\uparrow,\pi\downarrow$
\task $x\downarrow,\pi\downarrow$
\end{tasks}
\ExplainChoices{Wrong.}{Correct.}{Wrong.}{Wrong.}

\Question \textbf{Larger $\kappa$ means (holding $x$ fixed):}
\answer{B}
\begin{tasks}(1)
\task Inflation responds less
\task Inflation responds more
\task No change
\task Inflation becomes backward-looking
\end{tasks}
\ExplainChoices{Wrong.}{Correct.}{Wrong.}{Wrong.}

\vspace{0.6em}\Hrule\vspace{0.8em}
\section*{Problem 2 (Questions 21--40): NKPC + Taylor rule}
\Question \textbf{Numerical NKPC: $\beta=1$, $\mathbb{E}_0\pi_1=1\%$, $\kappa=0.5$, $x_0=2\%$, $u_0=0$. Then $\pi_0=$}
\answer{B}
\begin{tasks}(1)
\task $1\%$
\task $2\%$
\task $2\%$
\task $0.5\%$
\end{tasks}
\ExplainChoices{Wrong: forgot gap term.}{Correct: $1\%+0.5	imes2\%=2\%$.}{Duplicate value; treated as distractor.}{Wrong.}

\Question \textbf{Numerical Taylor: $ar i=1\%$, $\phi_\pi=1.5$, $\phi_x=0$, $\pi_0=2\%$. Then $i_0=$}
\answer{C}
\begin{tasks}(1)
\task $2\%$
\task $3\%$
\task $4\%$
\task $1\%$
\end{tasks}
\ExplainChoices{Wrong.}{Wrong.}{Correct.}{Wrong.}

\Question \textbf{Compute real rate: $i_0=4\%$, $\mathbb{E}_0\pi_1=1\%$. Then $r_0=$}
\answer{A}
\begin{tasks}(1)
\task $3\%$
\task $5\%$
\task $-3\%$
\task $1\%$
\end{tasks}
\ExplainChoices{Correct.}{Wrong.}{Wrong.}{Wrong.}

\Question \textbf{If CB raises $i_0$ holding expectations fixed, IS implies $x_0$:}
\answer{B}
\begin{tasks}(1)
\task Increases
\task Decreases
\task Unchanged
\task Undefined
\end{tasks}
\ExplainChoices{Wrong.}{Correct.}{Wrong.}{Wrong.}

\Question \textbf{If terminal conditions are $x_1=0$ and $\pi_1=0$, then:}
\answer{A}
\begin{tasks}(1)
\task $\mathbb{E}_0 x_1=\mathbb{E}_0\pi_1=0$
\task Expectations become larger
\task Expectations are indeterminate
\task Expectations equal current values
\end{tasks}
\ExplainChoices{Correct.}{Wrong.}{Wrong.}{Wrong.}

\Question \textbf{If $\phi_x$ increases, policy responds more to:}
\answer{B}
\begin{tasks}(1)
\task Inflation only
\task Output gap (stabilization)
\task Expected inflation only
\task Natural rate only
\end{tasks}
\ExplainChoices{Wrong.}{Correct.}{Wrong.}{Wrong.}

\Question \textbf{If $\phi_\pi=0$ (peg), higher inflation tends to:}
\answer{B}
\begin{tasks}(1)
\task Raise real rate
\task Lower real rate
\task Leave real rate unchanged
\task Stabilize inflation
\end{tasks}
\ExplainChoices{Wrong.}{Correct.}{Wrong.}{Wrong.}

\Question \textbf{With $\phi_\pi>1$, inflation rise tends to:}
\answer{B}
\begin{tasks}(1)
\task Reduce real rate
\task Increase real rate
\task Not change real rate
\task Set $x_0=0$ always
\end{tasks}
\ExplainChoices{Wrong.}{Correct.}{Wrong.}{Wrong.}

\Question \textbf{If expected inflation rises but rule reacts to current $\pi_0$ only, $i_0$ changes:}
\answer{B}
\begin{tasks}(1)
\task Immediately due to expectations
\task Only if $\pi_0$ changes
\task One-for-one always
\task Never
\end{tasks}
\ExplainChoices{Wrong.}{Correct.}{Wrong.}{Wrong.}

\Question \textbf{Which blocks are forward-looking?}
\answer{C}
\begin{tasks}(1)
\task IS only
\task NKPC only
\task Both IS and NKPC
\task Neither
\end{tasks}
\ExplainChoices{Wrong.}{Wrong.}{Correct.}{Wrong.}

\Question \textbf{If $u_0>0$ and CB wants lower $\pi_0$, it chooses:}
\answer{B}
\begin{tasks}(1)
\task $x_0>0$
\task $x_0<0$
\task $x_0=0$
\task Any $x_0$
\end{tasks}
\ExplainChoices{Wrong.}{Correct.}{Wrong.}{Wrong.}

\Question \textbf{If $\beta$ rises toward 1, NKPC becomes:}
\answer{B}
\begin{tasks}(1)
\task Less forward-looking
\task More forward-looking
\task Backward-looking
\task Unrelated to expectations
\end{tasks}
\ExplainChoices{Wrong.}{Correct.}{Wrong.}{Wrong.}

\Question \textbf{If $\mathbb{E}_0\pi_1$ rises holding $x_0$ fixed, $\pi_0$ rises by:}
\answer{B}
\begin{tasks}(1)
\task 0
\task $\beta$ times the change
\task $(1-\beta)$ times the change
\task $\kappa$ times the change
\end{tasks}
\ExplainChoices{Wrong.}{Correct.}{Wrong.}{Wrong.}

\Question \textbf{If $\kappa=0$, then inflation is:}
\answer{A}
\begin{tasks}(1)
\task Independent of $x_0$
\task More sensitive to $x_0$
\task Always equal to target
\task Always zero
\end{tasks}
\ExplainChoices{Correct.}{Wrong.}{Wrong.}{Wrong.}

\Question \textbf{If $\sigma$ is very large, IS response to real-rate gap is:}
\answer{B}
\begin{tasks}(1)
\task Strong
\task Weak
\task Zero always
\task Opposite sign
\end{tasks}
\ExplainChoices{Wrong.}{Correct.}{Wrong.}{Wrong.}

\Question \textbf{If $r_0^n$ increases holding $i_0$ and expectations fixed, then $x_0$:}
\answer{B}
\begin{tasks}(1)
\task Falls
\task Rises
\task Unchanged
\task Depends only on NKPC
\end{tasks}
\ExplainChoices{Wrong.}{Correct.}{Wrong.}{Wrong.}

\Question \textbf{If CB sets $i_0=r_0^n+\mathbb{E}_0\pi_1$ and $x_1=0$, then $x_0=$}
\answer{B}
\begin{tasks}(1)
\task $>0$
\task $=0$
\task $<0$
\task Undefined
\end{tasks}
\ExplainChoices{Wrong.}{Correct.}{Wrong.}{Wrong.}

\Question \textbf{If $i_0=0$ and $\mathbb{E}_0\pi_1=-2\%$, then $r_0=$}
\answer{C}
\begin{tasks}(1)
\task $-2\%$
\task $0\%$
\task $2\%$
\task $-4\%$
\end{tasks}
\ExplainChoices{Wrong.}{Wrong.}{Correct.}{Wrong.}

\Question \textbf{A higher $\phi_\pi$ tends to make equilibrium inflation (given shocks):}
\answer{B}
\begin{tasks}(1)
\task More sensitive
\task Less sensitive
\task Unchanged
\task Always zero
\end{tasks}
\ExplainChoices{Wrong.}{Correct.}{Wrong.}{Wrong.}

\Question \textbf{Anchored expectations roughly means:}
\answer{A}
\begin{tasks}(1)
\task $\mathbb{E}_0\pi_1$ stays near target
\task $\pi_0$ always zero
\task $x_0$ always zero
\task $i_0$ never changes
\end{tasks}
\ExplainChoices{Correct.}{Wrong.}{Wrong.}{Wrong.}

\end{Exercise}

\end{document}
